\begin{pastexam}[label={px:ntu-109-comm-dsa-46}]{臺大電信所 109 學年度（B）第 46 題}
（複選；每題 2 分，每答錯一個選項倒扣 0.4 分）

In a red-black tree of size n
\begin{enumerate}[label=(\Alph*),nosep]
\item Time complexity of one search is $O(\log n)$.
\item Time complexity of one insertion is $O(\log n)$.
\item Time complexity of n insertions is $O(n\log n)$.
\item Time complexity of one deletion is $O(\log n)$.
\item Time complexity of n deletions is $O(n\log n)$.
\end{enumerate}
\end{pastexam}

\begin{pastsoln}
答案：(A)(B)(C)(D)(E)，全對。

和第~44~題的伸展樹對照：紅黑樹保證的是\textbf{最壞情形}的高度。 $n$ 個內部節點的紅黑樹高度
不超過 $2\lg(n + 1)$（資結書定理~10.4），所以：
\begin{itemize}
\item (A) 搜尋從根走到葉，比較次數不超過高度，$O(\log n)$。
\item (B) 插入：往下找位置 $O(\log n)$，修正時重新著色沿路往上，最多 $O(\log n)$ 次，旋轉最多
$2$ 次（資結書第~10.3.4~節），合計 $O(\log n)$。
\item (C) 每次插入時樹最多 $2n$ 個節點，每次 $O(\log 2n) = O(\log n)$，$n$ 次 $O(n\log n)$。
\item (D) 刪除：找到節點、找後繼 $O(\log n)$，修正時往上最多 $O(\log n)$ 步，旋轉最多 $3$ 次
（資結書第~10.3.7~節），$O(\log n)$。
\item (E) 每次 $O(\log n)$，$n$ 次 $O(n\log n)$。
\end{itemize}
這些都是單一操作的最壞情形，不是平攤。
\end{pastsoln}
